\section{面積}\label{節:曲線間の面積}

\sectionref{節:積分符号面積}では,曲線と $x$ 軸の間の面積を定積分で表した.では,2つの曲線に挟まれた部分の面積はどう求めればよいだろうか.

$a<b$ とし,\,$f,g$ は $[a,b]$ で連続で,\,$f(x)\ge g(x)$ を満たすとする.求めるのは
\[
a\le x\le b,\qquad g(x)\le y\le f(x)
\]
で表される部分の面積 $S$ である.この部分を $x$ 軸に垂直な細い帯に分け,各帯を長方形で近似する.長方形の面積の和の極限として,上下の差 $f(x)-g(x)$ の定積分に帰着させることを目標とする.

区間 $[a,b]$ を $n$ 等分し,\,$\Delta x=\frac{b-a}{n}$,\,$x_k=a+k\Delta x$ と置く.\,$k$ 番目の帯は,幅 $\Delta x$,高さ $f(x_k)-g(x_k)$ の長方形で近似できる.よって
\[
S\approx\sum_{k=1}^n\{f(x_k)-g(x_k)\}\Delta x
\]
である.右辺は連続関数 $f-g$ のリーマン和なので,\,$n\to\infty$ として
\begin{equation}\label{式:上下の差の面積}
S=\int_a^b\{f(x)-g(x)\}\dd x
\end{equation}
を得る.

長方形の高さは上の値から下の値を引いた差であり,曲線が $x$ 軸の上にあるかどうかは使っていない.実際,\,$f,g$ に同じ定数 $c$ を足しても,\,$(f+c)-(g+c)=f-g$ は変わらない.したがって,上下の曲線がともに $x$ 軸の下にあっても,同じ式が使えるのだ.

ここで1つ断りをいれておこう.式\eqref{式:上下の差の面積}は,区間全体で $f(x)\ge g(x)$ であることを前提にしている.2曲線が途中で交差すると,上下が入れ替わった区間では $f-g$ が負になり,面積を引いてしまう.1つの式で書けば
\[
S=\int_a^b|f(x)-g(x)|\dd x
\]
であるが,実際に計算するには,上下関係が変わる点で区間を分けて絶対値を外す.交点の $x$ 座標は区切りの候補であり,交点を求めただけで上下が決まるわけではない.差を因数分解するか,各区間で1つの値を代入して符号を確かめる.

\begin{bookpoint}
\textbf{POINT}\par
交点の $x$ 座標で区間を分け,各区間で(上の関数)$-$(下の関数)を積分して足す.
\end{bookpoint}

\begin{example}{交差する曲線の面積}
2曲線 $y=x^3-3x$ と $y=x$ に囲まれた部分の面積を求めよ.
\end{example}

\paragraphbf{方針}
交点で区間を分け,差 $x^3-4x$ の符号を調べて絶対値を外す.

\paragraphbf{解答}
$8$

\paragraphbf{解法}
交点の $x$ 座標は
\[
x^3-3x=x
\]
より,
\[
x(x+2)(x-2)=0
\]
であるから,\,$x=-2,0,2$ である.差 $x^3-4x=x(x+2)(x-2)$ は $-2<x<0$ で正,\,$0<x<2$ で負である.よって面積は
\[
S=\int_{-2}^0(x^3-4x)\dd x+\int_0^2(4x-x^3)\dd x
\]
となる.図\ref{図:面積の上下交換}では,左右の部分が原点について対称である.積分でも,第1項で $x=-t$ と置けば
\[
\int_{-2}^0(x^3-4x)\dd x=\int_0^2(4t-t^3)\dd t
\]
となり,第2項と一致する.したがって
\begin{align*}
S&=2\int_0^2(4x-x^3)\dd x\\
 &=2\left[2x^2-\frac{x^4}{4}\right]_0^2\\
 &=2(8-4)\\
 &=8.
\end{align*}

\begin{figure}[htbp]
\centering
\begin{tikzpicture}[x=1.6cm,y=0.65cm,>=Stealth]
\fill[blue!12]plot[domain=-2:0,samples=60](\x,{\x*\x*\x-3*\x})--plot[domain=0:-2,samples=30](\x,\x)--cycle;
\fill[orange!17]plot[domain=0:2,samples=30](\x,\x)--plot[domain=2:0,samples=60](\x,{\x*\x*\x-3*\x})--cycle;
\draw[->](-2.3,0)--(2.4,0)node[right]{$x$};
\draw[->](0,-2.8)--(0,3.4)node[above]{$y$};
\draw[blue!70!black,thick]plot[domain=-2.13:2.13,samples=120](\x,{\x*\x*\x-3*\x});
\draw[orange!80!black,thick](-2.2,-2.2)--(2.2,2.2);
\foreach \x in {-2,2}{\draw[dashed](\x,0)--(\x,\x);\node[below]at(\x,0){$\x$};}
\node[below right]at(0,0){$O$};
\node[blue!70!black]at(-1.25,2.8){$y=x^3-3x$};
\node[orange!80!black]at(1.6,3.0){$y=x$};
\end{tikzpicture}
\caption{同じ2曲線でも,交点の左右で上側の曲線が入れ替わる.}\label{図:面積の上下交換}
\end{figure}

区間を分けずに $\int_{-2}^2(x^3-4x)\dd x$ とすると,奇関数の積分なので $0$ になる.この $0$ は,左の部分の面積を正,右の部分の面積を負として足した値である.面積を求めるには,符号が変わる $x=0$ で区間を分けなければならない.

\subsection{積分変数の選択}

細い帯は縦に取る必要はない.\,$c\leq y\leq d$ において,右側の境界が $x=q(y)$,左側が $x=p(y)$ と表せるなら,帯の横の長さは $q(y)-p(y)$ なので
\begin{equation}\label{式:左右の差の面積}
S=\int_c^d\{q(y)-p(y)\}\dd y
\end{equation}
である.式\eqref{式:上下の差の面積}では上下の差,この式では左右の差を使う.変数の文字を交換するだけでなく,実際にどの向きに帯を取るかが変わっている.

たとえば $x=y^2$ と $x=2-y$ に囲まれる部分では,交点の $y$ 座標は $-2,1$ である.その間では直線が右側なので,
\[
S=\int_{-2}^1(2-y-y^2)\dd y
 =\left[2y-\frac{y^2}{2}-\frac{y^3}{3}\right]_{-2}^1
 =\frac92.
\]
縦に切ると放物線の上側と下側を $y=\pm\sqrt x$ に分け,さらに直線との関係も場合分けすることになる.横に切れば,一つの区間で一つの差になる.境界をどちらの変数で表すと少ない場合分けで済むかが,積分変数を選ぶ基準になる.



%
